\documentclass[aspectratio=169,11pt]{beamer}
\usepackage{academia-slides}
\slidesmeta{Elementos de Philosophia}{Thiago Sinibaldi}{Simplicidade da alma humana}

\begin{document}

\slidecapa{Simplicidade da alma humana}{Psychologia, Livro I}{Roma, 1897 · 5.ª edição}

\slidesecao{Artigo II.}{Simplicidade da alma humana.}

\begin{slide}{tese}{Simplicidade da alma humana}{n. 29 · p. 35}
\begin{enunciado}
A alma humana é \textit{simples}, isto é, destituída de partes quantitativas e de composição.
\end{enunciado}

Esta tese dirige-se contra os \textsc{materialistas}, que reputam o pensamento um producto do cérebro.
\end{slide}

\begin{slide}{definicao}{Definição de substancia espiritual}{n. 30 · p. 36}
\begin{enunciado}
\textbf{Espiritual} diz-se a substancia que, tanto na existência como na operação, não depende intrinsecamente da matéria.
\end{enunciado}
Distingue-se do \textit{corpo}, cuja essência consiste na extensão.
\end{slide}

\begin{slide}{divisao}{Divisão do ente}{n. 31 · p. 37}
O ente divide-se:
\begin{topicos}
\item em \textit{substancia} e \textit{accidente}, segundo a existência em si ou em outro;
\item a substancia, por sua vez, em
  \begin{topicos}
  \item \textit{completa}, como o homem;
  \item \textit{incompleta}, como a alma separada;
  \end{topicos}
\item em \textit{acto} e \textit{potência}, segundo a perfeição.
\end{topicos}
\end{slide}

\begin{slide}{argumento}{Argumento}{n. 32 · p. 38}
\begin{rotulados}
\rot{Maior}{Todo ente composto de partes pode dissolver-se.}
\rot{Menor}{Ora a alma humana não é composta de partes.}
\rot{Conclusão}{Logo, a alma humana não pode dissolver-se.}
\end{rotulados}
\end{slide}

\begin{slide}{argumento}{Argumento (premissas longas)}{n. 32 · p. 38}
\begin{rotulados}
\rot{Maior}{O que é capaz de conhecer as essências universaes e immateriaes, e de reflectir sobre si mesmo em acto completo de reversão, não pode ser composto de partes extensas, nem depender intrinsecamente de órgão corporeo.}
\rot{Menor}{Ora a alma humana conhece as essências universaes e reflecte sobre si mesma plenamente.}
\rot{Conclusão}{Logo, a alma humana não é composta de partes extensas.}
\end{rotulados}
\end{slide}

\begin{slide}{objecao}{Objecção I}{n. 33 · p. 39}
O pensamento varia com o estado do cérebro; logo, é função d'elle. \textit{Mais: a alma adormece com o corpo.}
\end{slide}

\begin{slide}{resposta}{Resposta á objecção I}{n. 33 · p. 39}
Concedo o antecedente quanto ao \textit{exercício} da imaginação, nego quanto ao acto próprio do intellecto. A dependência é \textit{extrínseca}, não intrínseca.
\end{slide}

\begin{slide}{corolario}{Corollario}{n. 34 · p. 40}
Segue-se que a alma humana é \textit{incorruptível por si}, pois só se corrompe o que tem partes.
\end{slide}

\begin{slide}{escolio}{Escholio}{n. 34 · p. 40}
Os antigos, como Platão, chegaram a esta verdade por via diversa; S.~Thomás, porém, funda-a na \textit{operação}.
\end{slide}

\begin{slide}{observacao}{Observação}{n. 35 · p. 41}
Note-se que «simples» não equivale a «infinito»: a alma é simples e, todavia, finita.
\end{slide}

\begin{slide}{citacao}{S. Thomás}{Summa Theol., I, q. 75, a. 2}
\begin{citacao}
Anima humana, quæ dicitur intellectus vel mens, est quiddam incorporeum et subsistens. Intellectus enim cognoscit omnium corporum naturas; quod autem potest cognoscere aliqua, oportet ut nihil eorum habeat in sua natura.
\end{citacao}
\end{slide}

\begin{slide}{texto}{Texto corrido}{p. 42}
Paragrapho comum sem etiqueta. A questão da origem das almas foi muito agitada entre os antigos.

Segundo parágrapho, separado por linha em branco, para verificar o espaçamento entre blocos.
\end{slide}

\begin{slide}{tese}{Prova de que a alma humana, enquanto princípio de operações intellectuaes, é substancia simples e espiritual}{n. 36 · p. 43}
\begin{enunciado}
A alma humana é substancia espiritual.
\end{enunciado}
\begin{topicos}
\item Primeira prova, tirada da \textit{universalidade} das ideias intellectuaes;
\item segunda prova, tirada da \textit{reflexão} completa.
\end{topicos}
\end{slide}

\begin{slide}{texto}{Slide com muito texto}{n. 37 · p. 44}
A alma humana, considerada em si mesma, apresenta propriedades que nenhuma matéria pode explicar. Em primeiro logar, o intellecto abstrahe das condições individuaes e attinge o universal, o que exige princípio immaterial. Em segundo logar, a vontade tende ao bem universal, e não se satisfaz com nenhum bem particular. Em terceiro logar, a consciência reflecte sobre os próprios actos. Estas três propriedades, tomadas em conjunto, provam que o princípio de que dimanam é simples. Os materialistas respondem, porém, que tudo isto se reduz a movimentos moleculares; mas tal resposta confunde o movimento com o conhecimento do movimento.

Demais, a experiência quotidiana mostra que o homem conhece a si mesmo como um só e idêntico sujeito, através de todas as mudanças do corpo, o que seria impossível se a alma fosse mera collecção de partes em fluxo contínuo, como o corpo, que se renova inteiramente em poucos annos.
\continua
\end{slide}

\begin{slide}{divisao}{Divisão das opiniões}{n. 38 · p. 45}
\begin{rotulados}
\rot{a)}{Os \textit{materialistas}, que negam a espiritualidade.}
\rot{b)}{Os \textit{espiritualistas}, que a affirmam.}
\rot{c)}{Os \textit{scepticos}, que suspendem o juízo.}
\end{rotulados}
\begin{topicos}
\item Oito itens de teste de uma só linha, para medir o orçamento vertical dos tópicos aqui;
\item segundo item de teste de cerca de oitenta caracteres para medir a altura;
\item terceiro item de teste de cerca de oitenta caracteres para medir a altura;
\item quarto item de teste de cerca de oitenta caracteres para medir a altura;
\item quinto item de teste de cerca de oitenta caracteres para medir a altura;
\item sexto item de teste de cerca de oitenta caracteres para medir a altura;
\item sétimo item de teste de cerca de oitenta caracteres para medir a altura;
\item oitavo item de teste de cerca de oitenta caracteres para medir a altura.
\end{topicos}
\end{slide}

\end{document}
